\section{模型训练的评估困境}

\subsection{Benchmark 失效与 Agent 泛化}

\begin{warningbox}{Benchmark 过拟合的风险}
当前 Agent 领域面临严重的评估问题：
\begin{itemize}
    \item 可用的 Agent Benchmark 数量有限
    \item 刷高 Benchmark 分数不代表真实能力——很多时候是对特定任务的过拟合
    \item 用户在 OOD（Out-of-Distribution）场景的体感与 Benchmark 分数脱节
    \item "种瓜得瓜、种豆得豆"——训练什么就提升什么，但泛化有限
\end{itemize}
\end{warningbox}

\subsection{评估困境的根源}

杨植麟认为，评估困境的根源是：

\begin{enumerate}
    \item 强化学习最终依赖于好的评估/验证——数学题和有 Test Case 的编程可以自动验证，但复杂的端到端任务很难找到好的评估方式
    \item 模型能力提升带来的新问题：任务复杂度增加，评估方式需要跟上
    \item Agent 任务的长尾分布特性：难以覆盖所有真实场景
\end{enumerate}

\subsection{本章小结}

Benchmark 失效和 Agent 泛化不足是当前最突出的评估困境。根源在于复杂任务难以自动验证，以及训练分布与真实使用场景之间的 Gap。用 AI 参与评估和训练（AI 训练 AI）可能是突破方向。


